% State of the Art section. Paste into your own thesis, or \input it.
%
% Moved from clinical_background/06_geometry_risk.tex per docs_thesis/state_of_the_art_plan.md.
% \label{sec:geometry} kept alive -- referenced from clinical_background/00_chapter_intro.tex,
% clinical_background/03_cad.tex and clinical_background/05_hemodynamics.tex. The closing
% "Branching pattern" subsection of the original file was cut from here and expanded instead into
% 08_the_gap.tex, which is where docs_thesis/state_of_the_art_plan.md and find_research.md's own
% "The gap, stated for \S4" put that argument -- keeping it here would have said the same thing
% twice in adjacent sections.
%
% Verification status carried over unchanged: only shen2026geometry (cited in 08_the_gap.tex, not
% here after the split) and griffo2026wss (06_why_shape_enough.tex) are full-text read. Everything
% else below is [BIBLIOGRAPHY VERIFIED], abstract-level, per docs_thesis/find_research.md.
%
% Figure: figures/external/bifurcation_angle.png <- python scripts/fetch_external_figures.py
% Credit line copied verbatim from figures/external/CREDITS.md -- never typed by hand.
%
% STILL WANTED, ours not external: the four measurements drawn on one tree
% (scripts/make_tree_figure.py, part of the topology rebuild).

\section{Does tree geometry predict disease?}
\label{sec:geometry}

The shape of the coronary tree is a clinical quantity, not just a description of it.
Shape sets the drag on the wall, the drag decides how the lining behaves, and the lining is what
produces plaque (Section~\ref{sec:hemodynamics}) \cite{rampidis2022geometry}.
What makes a particular measurement worth taking is evidence that the chain holds in patients.
Three measurements have that evidence, and none of them has as much of it as a first reading
suggests.

\subsection{Evidence from patients}

Han et al.\ looked at patients who went on to have a heart attack and asked what about the culprit
lesion predicted it, beyond how narrow the vessel was and what the plaque was made of
\cite{han2022plaque}.
Three things did: the lesion sitting closer to the ostium, the lesion sitting at a bifurcation, and
higher tortuosity.
None of the three is a property of the plaque; all three describe the tree it grew in, which is
why the shape of the tree is worth studying on its own.

Stone et al.\ then showed the shape comes first.
They measured shear inside the arteries of 506 patients and reimaged 374 of them six to ten months
later \cite{stone2012prediction}.
Low shear at baseline predicted which segments went on to lose lumen area.
Two limits matter: it did not predict the primary endpoint, growth in plaque area, which went with
how much plaque was already there; and every patient had already had a heart attack and been
stented, so this does not transfer unchanged to people who are still well.

\subsection{Bifurcation angle}

The left main bifurcation angle has the firmest evidence of the three, from five separate cohorts
pointing the same way.
Temov and Sun measured it on CCTA in 196 patients and found a mean of
$79.40^\circ \pm 22.97^\circ$, spanning $35.5^\circ$ to $178^\circ$ \cite{temov2016bifurcation}.
The angle varies enormously between healthy people, and men were twice as likely to exceed
$80^\circ$, so any distribution reported from a cohort has to be split by sex rather than pooled
\cite{temov2016bifurcation}.

The angle widens with disease.
Juan et al.\ sorted 313 CCTA patients by how much disease they had and found
$75.5^\circ \pm 14.7^\circ$ in the 40 with clean arteries, $81.2^\circ \pm 20.1^\circ$ in the 62
whose plaque did not block the vessel, and $87.3^\circ \pm 18.8^\circ$ in the 211 with a stenosis of
at least half the vessel \cite{juan2017bifurcation}.
Mansouri et al.\ replicated the gradient against a continuous severity score in 122 patients who had
both CCTA and invasive angiography: the angle correlated with the Gensini score at 0.684, ahead of
plaque volume at 0.281, remodeling index at 0.438 and calcium score at 0.217
\cite{mansouri2024bifurcation}.
A geometric measurement outperforming the plaque measurements in the same patients is the strongest
argument in this literature for measuring shape at all.

One study reaches past disease severity to events.
Tommasino et al.\ followed 499 consecutive CCTA patients and found a left main bifurcation angle
above $80^\circ$ to be the strongest predictor of a major adverse cardiac event, ahead of diabetes,
obstructive disease and high-risk plaque \cite{tommasino2024clap}.
Wang et al.\ supply the mechanism: they simulated the flow through 76 left main arteries sorted into
four shapes, and shear was lowest around the branch points of the LAD and lowest of all in the
wide-angled shape \cite{wang2024lmphenotypes}.
An anatomical gradient, an outcome association and a hemodynamic mechanism pointing the same way are
what make this angle worth measuring in its own right.

On its own the angle discriminates only modestly. Cui et al.\ measured 318 bifurcation angles in
106 patients and found it an independent predictor of significant left coronary stenosis, at 67\%
sensitivity and 78\% specificity \cite{cui2017bifurcation} --- a population descriptor rather than a
diagnostic test, and treated as one here.

\begin{figure}[h]
  \centering
  \includegraphics[width=0.85\textwidth]{content/background/figures/bifurcation_angle.png}
  \caption{Two vessels can share a bifurcation angle and disagree about what it is. Panels A and B
    measure the same junction by different conventions --- $\alpha$ from the extension of the
    proximal main vessel, $\beta$ from the distal main vessel --- and panel C shows the two
    diverging once the main vessel curves. MV: main vessel; SB: side branch. Figure from
    Murasato Y, Meno K, Mori T, Tanenaka K, PLOS ONE 2022, doi:10.1371/journal.pone.0273157,
    CC BY 4.0.}
  \label{fig:bifurcation-angle}
\end{figure}

\subsection{Tortuosity}

Tortuosity is the second measurement, and its association with disease runs in opposite directions
depending on which disease is meant.
Groves et al.\ reviewed 1221 consecutive angiograms, the largest tortuosity cohort available, and
found severe tortuosity in 12.45\% of them --- associated with a significantly \emph{lower} rate of
significant coronary disease \cite{groves2009tortuosity}.
Zebi\'c Mihi\'c et al.\ studied 131 patients being investigated for suspected disease and found the
opposite for disease that does \emph{not} obstruct the vessel, at roughly eight times the odds, an
odds ratio of 7.96 with a 95\% confidence interval from 1.79 to 35.40
\cite{zebicmihic2023tortuosity}.
The two agree once the endpoint is named precisely: tortuosity tracks ischemia without obstruction,
and tracks inversely with obstruction itself. Reporting only one would mislead every tortuosity
number this thesis goes on to report.

How tortuosity is measured decides what is found.
A second study from the same group measured it two ways in the same 160 patients, by a continuous
index and by the clinical count of bends sharper than $45^\circ$ \cite{zebicmihic2024index}.
The continuous index localized to the territory that was ischemic --- highest in the circumflex with
lateral ischemia, highest in the LAD with anterior ischemia --- while the bend count found almost
nothing.
This thesis therefore computes a continuous index on the centerlines rather than reproducing the
clinical bend score.

\subsection{Dominance}

Dominance is the third, and the evidence for it is weaker and more conditional than the most-cited
study suggests.
Veltman et al.\ followed 1425 patients after CCTA for a median of two years and reported that
left-dominant patients reached a heart attack or death about three times as often as right-dominant
ones \cite{veltman2012dominance}.
The CONFIRM registry does not reproduce it.
Following 6382 CCTA patients across twelve centers for five years, survival did not differ by
dominance in patients with obstructive disease, without it, or with normal arteries
\cite{gebhard2015dominance}.
The one positive result was a subgroup: left dominance together with left main disease.

A meta-analysis of 255,718 patients with acute coronary syndrome finds a real but small effect,
odds ratio 1.27 for mortality \cite{khan2016dominance}.
Read together, the three say dominance matters once an infarct has happened --- a left-dominant
heart has more muscle depending on the left system --- and not as a predictor of getting there. It is
used here to stratify the other features, not reported as a risk factor.

\subsection{Two limits on all three}

The first is that the measurement error is the size of the effect.
Givehchi et al.\ imaged nine phantoms with known bifurcation angles and recovered them to within a
few degrees, then applied the same two accepted techniques to 50 real CCTA cases, where they
disagreed with each other by $12.0^\circ \pm 10.6^\circ$ \cite{givehchi2018phantom}.
Juan's entire gradient from clean arteries to significant stenosis spans about $12^\circ$.
A published angle--disease association is therefore only as real as the protocol behind it, and any
angle reported here must come from a single automatic definition computed on the centerline rather
than from a manual reformation.
It is also why measuring our own reproducibility is a prerequisite rather than an extra.

The second is that the geometry is not fixed while the disease develops.
Glagov et al.\ examined the left main artery in 136 hearts at autopsy and found that arteries enlarge
as plaque grows: the lumen does not narrow at all until the lesion occupies about 40\% of the vessel,
and only then falls sharply \cite{glagov1987remodeling}.
Every centerline and surface in this thesis is derived from a lumen segmentation, so early disease is
nearly invisible to it while advanced disease has already reshaped the very geometry being measured.
A cross-sectional association between shape and disease is therefore partly the disease writing its
own predictor.
